\textbf{(i)} Both ideals consist of finite sums $\sum a_ix^i$ with $a_i\in\mathfrak{a}$.

\textbf{(ii)} We have $A[x]/\mathfrak{p}[x]\cong(A/\mathfrak{p})[x]$, which is an integral domain.

\textbf{(iii)} Suppose $f$ is a zero-divisor in $(A/\mathfrak{q})[x]$. By \exref{1.2}, some nonzero $a\in A/\mathfrak{q}$ annihilates $f$. Every coefficient of $f$ is therefore a zero-divisor in $A/\mathfrak{q}$ and hence nilpotent. Again by \exref{1.2}, $f$ is nilpotent, so $\mathfrak{q}[x]$ is primary. The same coefficient criterion gives $r(\mathfrak{q}[x])=r(\mathfrak{q})[x]=\mathfrak{p}[x]$.

\textbf{(iv)} Comparing coefficients gives $(\bigcap_i\mathfrak{q}_i)[x]=\bigcap_i\mathfrak{q}_i[x]$. Distinct radicals remain distinct on extension, since contraction recovers the original ideal. For each $i$, a constant polynomial in $\bigcap_{j\ne i}\mathfrak{q}_j\setminus\mathfrak{q}_i$ proves irredundancy.

\textbf{(v)} If $\mathfrak{a}[x]\subseteq\mathfrak{q}\subseteq\mathfrak{p}[x]$ and $\mathfrak{q}$ is prime, then $\mathfrak{a}\subseteq\mathfrak{q}\cap A\subseteq\mathfrak{p}$. Minimality gives $\mathfrak{q}\cap A=\mathfrak{p}$, hence $\mathfrak{p}[x]\subseteq\mathfrak{q}$.
